\chapter{2015年陕西卷（文）}

\section{单选题}

\begin{problem}
设集合 \(M = \{x \mid x^2 = x\}\)，\(N = \{x \mid \lg x \leqslant 0\}\)，则 \(M \cup N =\)（\quad）

\choices
    {\([0,1]\)}
    {\((0,1]\)}
    {\([0,1)\)}
    {\((-\infty,1]\)}
\end{problem}

\begin{answer}
A
\end{answer}

\begin{solution}
方程 \(x^2=x\) 等价于 \(x(x-1)=0\)，所以 \(M=\{0,1\}\). 由于对数的真数必须为正，且 \(\lg x\leqslant0\) 等价于 \(0<x\leqslant1\)，所以 \(N=(0,1]\). 因此
\[
M\cup N=\{0,1\}\cup(0,1]=[0,1]
\]
故选 A.
\end{solution}

\begin{problem}
某中学初中部共有 \(110\text{ 名}\)教师，高中部共有 \(150\text{ 名}\)教师，其性别比例如图所示，则该校女教师的人数为（\quad）

\bitmapfigure[max width=0.56\linewidth]{img/2015/shaanxi_liberal/q2_fig1.png}

\choices
    {\(93\)}
    {\(123\)}
    {\(137\)}
    {\(167\)}
\end{problem}

\begin{answer}
C
\end{answer}

\begin{solution}
初中部女教师有 \(110\times70\%=77\) 人，高中部的扇形图中男教师占 \(60\%\)，所以女教师占 \(40\%\)，有 \(150\times40\%=60\) 人. 因此全校女教师共有 \(77+60=137\) 人，故选 C.
\end{solution}

\begin{problem}
已知抛物线 \(y^{2}=2px\)（\(p>0\)）的准线经过点 \((-1,1)\)，则该抛物线焦点坐标为（\quad）

\choices
    {\((-1,0)\)}
    {\((1,0)\)}
    {\((0,-1)\)}
    {\((0,1)\)}
\end{problem}

\begin{answer}
B
\end{answer}

\begin{solution}
抛物线 \(y^2=2px\) 的焦点为 \(\left(\frac{p}{2},0\right)\)，准线方程为 \(x=-\frac{p}{2}\). 因为点 \((-1,1)\) 在准线上，所以 \(-\frac{p}{2}=-1\)，解得 \(p=2\). 故焦点坐标为 \((1,0)\)，选 B.
\end{solution}

\begin{problem}
设 \(f(x) = \begin{cases}
1 - \sqrt{x}, & x \geqslant 0,\\
2^x, & x < 0,
\end{cases}\) 则 \(f(f(-2)) =\)（\quad）

\choices
    {\(-1\)}
    {\(\frac{1}{4}\)}
    {\(\frac{1}{2}\)}
    {\(\frac{3}{2}\)}
\end{problem}

\begin{answer}
C
\end{answer}

\begin{solution}
因为 \(-2<0\)，所以 \(f(-2)=2^{-2}=\frac{1}{4}\). 又因为 \(\frac{1}{4}\geqslant0\)，所以
\[
f\left(f(-2)\right)=f\left(\frac{1}{4}\right)=1-\sqrt{\frac{1}{4}}=1-\frac{1}{2}=\frac{1}{2}
\]
故选 C.
\end{solution}

\begin{problem}
一个几何体的三视图如图所示，则该几何体的表面积为（\quad）

\bitmapfigure[max width=0.42\linewidth]{img/2015/shaanxi_liberal/q5_fig1.png}

\choices
    {\(3\pi\)}
    {\(4\pi\)}
    {\(2\pi + 4\)}
    {\(3\pi + 4\)}
\end{problem}

\begin{answer}
D
\end{answer}

\begin{solution}
由三视图可知，该几何体是底面半径为 \(1\)、柱高为 \(2\) 的半圆柱. 其曲面展开后是一个长为半圆弧长 \(\pi\)、宽为 \(2\) 的矩形，面积为 \(2\pi\)；沿直径的矩形底面面积为 \(2\times2=4\)；两个半圆形端面的总面积为一个半径为 \(1\) 的圆的面积 \(\pi\). 因此表面积为
\[
2\pi+4+\pi=3\pi+4
\]
故选 D.
\end{solution}

\begin{problem}
“\(\sin \alpha = \cos \alpha\)”是“\(\cos 2\alpha = 0\)”的（\quad）

\choices
    {充分不必要条件}
    {必要不充分条件}
    {充分必要条件}
    {既不充分也不必要}
\end{problem}

\begin{answer}
A
\end{answer}

\begin{solution}
若 \(\sin\alpha=\cos\alpha\)，则 \(\sin^2\alpha=\cos^2\alpha\)，从而 \(\cos2\alpha=\cos^2\alpha-\sin^2\alpha=0\)，所以前者是后者的充分条件. 但取 \(\alpha=\frac{3\pi}{4}\)，有 \(\cos2\alpha=\cos\frac{3\pi}{2}=0\)，而 \(\sin\alpha\ne\cos\alpha\)，所以前者不是后者的必要条件. 故选 A.
\end{solution}

\begin{problem}
根据框图，当输入 \(x\) 为 \(6\) 时，输出的 \(y =\)（\quad）

\bitmapfigure[max width=0.26\linewidth]{img/2015/shaanxi_liberal/q7_fig1.png}

\choices
    {\(1\)}
    {\(2\)}
    {\(5\)}
    {\(10\)}
\end{problem}

\begin{answer}
D
\end{answer}

\begin{solution}
输入 \(x=6\) 后，先执行 \(x=x-3\)，得到 \(x=3\). 此时 \(x\geqslant0\)，沿“是”分支执行 \(y=x^2+1\)，所以 \(y=3^2+1=10\). 故选 D.
\end{solution}

\begin{problem}
对任意向量 \(\bs{a}\)，\(\bs{b}\)，下列关系式中不恒成立的是（\quad）

\choices
    {\( |\bs{a} \cdot \bs{b}| \leqslant |\bs{a}||\bs{b}| \)}
    {\( |\bs{a} - \bs{b}| \leqslant ||\bs{a}| - |\bs{b}|| \)}
    {\((\bs{a} + \bs{b})^2 = |\bs{a} + \bs{b}|^2\)}
    {\((\bs{a} + \bs{b})\cdot (\bs{a} - \bs{b}) = \bs{a}^2 - \bs{b}^2\)}
\end{problem}

\begin{answer}
B
\end{answer}

\begin{solution}
选项 A 是柯西不等式，恒有 \(\lvert\bs{a}\cdot\bs{b}\rvert\leqslant\lvert\bs{a}\rvert\lvert\bs{b}\rvert\). 选项 C 是向量平方的恒等式，选项 D 由分配律可得
\[
(\bs{a}+\bs{b})\cdot(\bs{a}-\bs{b})=\bs{a}\cdot\bs{a}-\bs{b}\cdot\bs{b}=\bs{a}^{2}-\bs{b}^{2}
\]
而选项 B 不恒成立. 例如取 \(\bs{a}=(1,0)\)，\(\bs{b}=(0,1)\)，则
\[
\lvert\bs{a}-\bs{b}\rvert=\sqrt{2}>0=\left\lvert\lvert\bs{a}\rvert-\lvert\bs{b}\rvert\right\rvert
\]
故选 B.
\end{solution}

\begin{problem}
设 \(f(x) = x - \sin x\)，则 \(f(x)\)（\quad）

\choices
    {既是奇函数又是减函数}
    {既是奇函数又是增函数}
    {是有零点的减函数}
    {是没有零点的奇函数}
\end{problem}

\begin{answer}
B
\end{answer}

\begin{solution}
因为 \(f(-x)=-x+\sin x=-f(x)\)，所以 \(f(x)\) 是奇函数. 又
\[
f'(x)=1-\cos x\geqslant0
\]
对任意 \(x_1<x_2\)，函数 \(1-\cos x\) 在区间 \([x_1,x_2]\) 上不恒为零，因此
\[
f(x_2)-f(x_1)=\int_{x_1}^{x_2}(1-\cos x)\,\mathrm{d}x>0
\]
所以 \(f(x)\) 是增函数. 且 \(f(0)=0\)，故选 B.
\end{solution}

\begin{problem}
设 \(f(x) = \ln x\)，\(0 < a < b\). 若 \(p = f\left(\sqrt{ab}\right)\)，\(q = f\left(\frac{a + b}{2}\right)\)，\(r = \frac{1}{2}(f(a) + f(b))\)，则下列关系式中正确的是（\quad）

\choices
    {\(q = r < p\)}
    {\(q = r > p\)}
    {\(p = r < q\)}
    {\(p = r > q\)}
\end{problem}

\begin{answer}
C
\end{answer}

\begin{solution}
由对数运算，
\[
p=\ln\sqrt{ab}=\frac{1}{2}(\ln a+\ln b)=r
\]
因为 \(0<a<b\)，由算术平均数大于几何平均数，得 \(\dfrac{a+b}{2}>\sqrt{ab}\). 函数 \(\ln x\) 在 \((0,+\infty)\) 上单调递增，所以 \(q>p\). 因此 \(p=r<q\)，选 C.
\end{solution}

\begin{problem}
某企业生产甲，乙两种产品均需用 A，B 两种原料，已知生产 \(1\text{ 吨}\)每种产品所需原料及每天原料的可用限额如表所示.

如果生产 \(1\text{ 吨}\)甲，乙产品可获利润分别为 \(3\text{ 万元}\)，\(4\text{ 万元}\)，则该企业每天可获得最大利润为（\quad）

\begin{center}
\begin{tabular}{|c|c|c|c|}
\hline
 & 甲 & 乙 & 原料限额 \\
\hline
 A（吨） & \(3\) & \(2\) & \(12\) \\
\hline
 B（吨） & \(1\) & \(2\) & \(8\) \\
\hline
\end{tabular}
\end{center}

\choices
    {\(12\text{ 万元}\)}
    {\(16\text{ 万元}\)}
    {\(17\text{ 万元}\)}
    {\(18\text{ 万元}\)}
\end{problem}

\begin{answer}
D
\end{answer}

\begin{solution}
设每天生产甲、乙两种产品分别为 \(x\) 吨、\(y\) 吨，其中 \(x\geqslant0\)，\(y\geqslant0\). 由原料限额得
\(3x+2y\leqslant12\)，\(x+2y\leqslant8\).
利润函数为 \(z=3x+4y\). 可行域的顶点为 \((0,0)\)、\((4,0)\)、\((0,4)\) 以及两条限额直线的交点. 联立
\(3x+2y=12\)，\(x+2y=8\)
得 \(x=2\)，\(y=3\). 各顶点处的利润分别为 \(0\)、\(12\)、\(16\)、\(18\) 万元，所以最大利润为 \(18\) 万元. 故选 D.
\end{solution}

\begin{problem}
设复数 \(z = (x - 1) + yi\)（\(x\)，\(y \in \R\)），若 \(|z| \leqslant 1\)，则 \(y \geqslant x\) 的概率为（\quad）

\choices
    {\(\frac{3}{4} + \frac{1}{2\pi}\)}
    {\(\frac{1}{2} + \frac{1}{\pi}\)}
    {\(\frac{1}{4} - \frac{1}{2\pi}\)}
    {\(\frac{1}{2} - \frac{1}{\pi}\)}
\end{problem}

\begin{answer}
C
\end{answer}

\begin{solution}
条件 \(|z|\leqslant1\) 等价于圆域
\[
(x-1)^2+y^2\leqslant1
\]
所求事件 \(y\geqslant x\) 对应圆心 \((1,0)\) 上方由直线 \(y=x\) 截出的圆弓形. 该直线到圆心的距离为 \(\dfrac{1}{\sqrt{2}}\)，弦两端为 \((0,0)\)，\((1,1)\)，对应的圆心角为 \(\dfrac{\pi}{2}\). 因此圆弓形面积为
\[
\frac{1}{2}\cdot1^2\cdot\frac{\pi}{2}-\frac{1}{2}\cdot1\cdot1=\frac{\pi}{4}-\frac{1}{2}
\]
总圆域面积为 \(\pi\)，故所求概率为
\[
\frac{\frac{\pi}{4}-\frac{1}{2}}{\pi}=\frac{1}{4}-\frac{1}{2\pi}
\]
故选 C.
\end{solution}

\section{填空题}

\begin{problem}
中位数为 \(1010\) 的一组数构成等差数列，其末项为 \(2015\)，则该数列的首项为\fillinblank{}.
\end{problem}

\begin{answer}
\(5\)
\end{answer}

\begin{solution}
设该等差数列共有 \(n\) 项，首项为 \(a_1\)，末项为 \(a_n=2015\). 等差数列的各项关于首、末两项的平均数成对对称，其中位数为
\[
\frac{a_1+a_n}{2}=1010
\]
因此 \(a_1+2015=2020\)，解得 \(a_1=5\).
\end{solution}

\begin{problem}
如图，某港口一天 \(6\text{ 时}\)到 \(18\text{ 时}\)的水深变化曲线近似满足函数 \(y = 3\sin \left(\frac{\pi}{6} x + \varphi\right) + k\).

据此函数可知，这段时间水深（单位：\(\mathrm{m}\)）的最大值为\fillinblank{}.

\begin{center}
\begin{minipage}{0.42\linewidth}
\centering
\bitmapfigure{img/2015/shaanxi_liberal/q14_fig1.png}
\end{minipage}
\end{center}
\end{problem}

\begin{answer}
\(8\)
\end{answer}

\begin{solution}
函数 \(y=3\sin\left(\frac{\pi}{6}x+\varphi\right)+k\) 的振幅为 \(3\)，所以最大值为 \(k+3\)，最小值为 \(k-3\). 由图可知这段时间水深的最小值为 \(2\)，故 \(k-3=2\)，从而 \(k=5\). 因此这段时间水深的最大值为 \(k+3=8\).
\end{solution}

\begin{problem}
函数 \(y = x\e^{x}\) 在其极值点处的切线方程为\fillinblank{}.
\end{problem}

\begin{answer}
\(y=-\frac{1}{\e}\)
\end{answer}

\begin{solution}
函数的导数为
\[
y'=\e^x(x+1)
\]
因为 \(\e^x>0\)，当 \(x<-1\) 时 \(y'<0\)，当 \(x>-1\) 时 \(y'>0\)，所以函数在 \(\left(-\infty,-1\right)\) 上递减、在 \(\left(-1,+\infty\right)\) 上递增，故 \(x=-1\) 是唯一极值点. 此时函数值为 \(y=-\frac{1}{\e}\)，且切线斜率 \(y'(-1)=0\). 因而切线方程为
\[
y=-\frac{1}{\e}
\]
\end{solution}

\begin{problem}
观察下列等式：

\(\displaystyle 1 - \frac{1}{2} = \frac{1}{2}\)；\(\displaystyle 1 - \frac{1}{2} + \frac{1}{3} - \frac{1}{4} = \frac{1}{3} + \frac{1}{4}\)；\(\displaystyle 1 - \frac{1}{2} + \frac{1}{3} - \frac{1}{4} + \frac{1}{5} - \frac{1}{6} = \frac{1}{4} + \frac{1}{5} + \frac{1}{6}\)，\(\cdots\)

据此规律，第 \(n\) 个等式可为\fillinblank{}.
\end{problem}

\begin{answer}
\(\displaystyle 1 - \frac{1}{2} + \frac{1}{3} - \frac{1}{4} + \cdots + \frac{1}{2n-1} - \frac{1}{2n} = \frac{1}{n+1} + \frac{1}{n+2} + \cdots + \frac{1}{2n}\)
\end{answer}

\begin{solution}
第 \(n\) 个等式左边为
\[
\sum_{j=1}^{n}\left(\frac{1}{2j-1}-\frac{1}{2j}\right)=\sum_{j=1}^{2n}\frac{1}{j}-2\sum_{j=1}^{n}\frac{1}{2j}=\sum_{j=1}^{2n}\frac{1}{j}-\sum_{j=1}^{n}\frac{1}{j}=\sum_{j=n+1}^{2n}\frac{1}{j}
\]
所以第 \(n\) 个等式为
\[
1-\frac{1}{2}+\frac{1}{3}-\frac{1}{4}+\cdots+\frac{1}{2n-1}-\frac{1}{2n}=\frac{1}{n+1}+\frac{1}{n+2}+\cdots+\frac{1}{2n}
\]
\end{solution}

\section{解答题}

\begin{problem}
\(\triangle ABC\) 的内角 \(A\)，\(B\)，\(C\) 所对的边分别为 \(a\)，\(b\)，\(c\). 向量 \(\bs{m} = (a, \sqrt{3}b)\) 与 \(\bs{n} = (\cos A, \sin B)\) 平行.
\begin{enumerate}
\item 求 \(A\)；
\item 若 \(a = \sqrt{7}\)，\(b = 2\)，求 \(\triangle ABC\) 的面积.
\end{enumerate}
\end{problem}

\begin{solution}
\begin{enumerate}
\item
由 \(\bs{m}\) 与 \(\bs{n}\) 平行，得
\[
a\sin B=\sqrt{3}b\cos A
\]
由正弦定理，\(\dfrac{a}{b}=\dfrac{\sin A}{\sin B}\)，所以 \(\sin A=\sqrt{3}\cos A\). 由于 \(0<A<\pi\)，有 \(\sin A>0\)，从而 \(\cos A>0\)，且 \(\tan A=\sqrt{3}\)，故 \(A=\dfrac{\pi}{3}\).

\item
当 \(a=\sqrt{7}\)，\(b=2\) 时，由余弦定理，
\[
a^2=b^2+c^2-2bc\cos A=4+c^2-2c
\]
因此 \(c^2-2c-3=0\)，解得 \(c=3\)（舍去负值）. 所以
\[
S_{\triangle ABC}=\frac{1}{2}bc\sin A=\frac{1}{2}\cdot2\cdot3\cdot\frac{\sqrt{3}}{2}=\frac{3\sqrt{3}}{2}
\]
\end{enumerate}
\end{solution}

\begin{problem}
如图 1，在直角梯形 \(ABCD\) 中，\(AD \parallel BC\)，\(\angle BAD = \frac{\pi}{2}\)，\(AB = BC = \frac{1}{2}AD = a\)，\(E\) 是 \(AD\) 的中点，\(O\) 是 \(AC\) 与 \(BE\) 的交点. 将 \(\triangle ABE\) 沿 \(BE\) 折起到图 2 中 \(\triangle A_{1}BE\) 的位置，得到四棱锥 \(A_{1}-BCDE\).

\begin{enumerate}
\item 证明：\(CD \perp\) 平面 \(A_{1}OC\)；
\item 若平面 \(A_{1}BE \perp\) 平面 \(BCDE\)，四棱锥 \(A_{1}-BCDE\) 的体积为 \(36\sqrt{2}\)，求 \(a\) 的值.
\end{enumerate}

\examfiguregroup[float=inline]{img/2015/shaanxi_liberal/q18_fig1.png}{%
    \makebox[0.44\linewidth][c]{\bitmapinclude[max width=\linewidth]{img/2015/shaanxi_liberal/q18_fig1.png}}\hfill
 \makebox[0.44\linewidth][c]{\bitmapinclude[max width=\linewidth]{img/2015/shaanxi_liberal/q18_fig2.png}}%
 }
 \end{problem}

\begin{solution}
\begin{enumerate}
\item
在折叠前的平面内建立直角坐标系，取
\(A=(0,0,0)\)，\(B=(0,-a,0)\)，\(C=(a,-a,0)\)，\(D=(2a,0,0)\)，\(E=(a,0,0)\).
于是 \(\overrightarrow{BE}=(a,a,0)\)，\(\overrightarrow{AC}=(a,-a,0)\)，从而 \(AC\perp BE\). 因为 \(O\) 在 \(AC\) 与 \(BE\) 上，所以 \(OC\perp BE\). 折叠是绕 \(BE\) 进行的，点 \(O\) 在折痕上，且折叠保持直角，因此 \(A_{1}O\perp BE\). 直线 \(A_{1}O\) 与 \(OC\) 在点 \(O\) 相交且都垂直于 \(BE\)，故 \(BE\perp\) 平面 \(A_{1}OC\). 又 \(CD\parallel BE\)，所以 \(CD\perp\) 平面 \(A_{1}OC\).

\item
四边形 \(BCDE\) 是平行四边形，底边 \(BC=a\)，高为 \(a\)，故其面积为 \(a^2\). 当平面 \(A_{1}BE\perp\) 平面 \(BCDE\) 时，四棱锥的高等于平面 \(A_{1}BE\) 内点 \(A_{1}\) 到 \(BE\) 的距离. 三角形 \(ABE\) 为直角等腰三角形，且 \(AB=AE=a\)，\(BE=a\sqrt{2}\)，所以该距离为
\[
h=\frac{2S_{\triangle ABE}}{BE}=\frac{a^2}{a\sqrt{2}}=\frac{a}{\sqrt{2}}
\]
于是四棱锥体积
\[
V=\frac{1}{3}a^2h=\frac{a^3}{3\sqrt{2}}=36\sqrt{2}
\]
解得 \(a^3=216\)，所以 \(a=6\).
\end{enumerate}
\end{solution}

 \begin{problem}
随机抽取一个年份，对西安市该年 \(4\) 月份的天气情况进行统计，结果如下：

\begingroup
\setlength{\tabcolsep}{3pt}
\begin{center}
\begin{tabular}{|c|c|c|c|c|c|c|c|c|c|c|}
\hline
日期 & \(1\) & \(2\) & \(3\) & \(4\) & \(5\) & \(6\) & \(7\) & \(8\) & \(9\) & \(10\) \\
\hline
天气 & 晴 & 雨 & 阴 & 阴 & 阴 & 雨 & 阴 & 晴 & 晴 & 晴 \\
\hline
日期 & \(11\) & \(12\) & \(13\) & \(14\) & \(15\) & \(16\) & \(17\) & \(18\) & \(19\) & \(20\) \\
\hline
天气 & 阴 & 晴 & 晴 & 晴 & 晴 & 晴 & 阴 & 雨 & 阴 & 阴 \\
\hline
日期 & \(21\) & \(22\) & \(23\) & \(24\) & \(25\) & \(26\) & \(27\) & \(28\) & \(29\) & \(30\) \\
\hline
天气 & 晴 & 阴 & 晴 & 晴 & 晴 & 阴 & 晴 & 晴 & 晴 & 雨 \\
\hline
\end{tabular}
\end{center}
\endgroup

\begin{enumerate}
\item 在 \(4\) 月份任取一天，估计西安市在该天不下雨的概率；
\item 西安市某学校拟从 \(4\) 月份的一个晴天开始举行连续 \(2\) 天的运动会，估计运动会期间不下雨的概率.
 \end{enumerate}
 \end{problem}

\begin{solution}
\begin{enumerate}
\item 表中 \(30\) 天中有 \(4\) 天下雨，因此任取一天不下雨的频率为 \(\dfrac{30-4}{30}=\dfrac{13}{15}\).
\item 表中晴天共有 \(16\) 天. 以晴天作为运动会第一天时，若第一天是 \(1\) 日或 \(29\) 日，则第二天分别为雨天 \(2\) 日或 \(30\) 日；其余 \(14\) 个晴天作为第一天时，连续两天均不下雨. 因此运动会期间不下雨的频率为 \(\dfrac{14}{16}=\dfrac{7}{8}\).
\end{enumerate}
\end{solution}

 \begin{problem}
 如图，椭圆 \(E: \frac{x^2}{a^2} + \frac{y^2}{b^2} = 1\) （\(a > b > 0\)）经过点 \(A(0,-1)\)，且离心率为 \(\frac{\sqrt{2}}{2}\).

\begin{enumerate}
\item 求椭圆 \(E\) 的方程；
\item 经过点 \((1,1)\)，且斜率为 \(k\) 的直线与椭圆 \(E\) 交于不同的两点 \(P\)，\(Q\)（均异于点 \(A\)），证明：直线 \(AP\) 与 \(AQ\) 的斜率之和为 2.
\end{enumerate}

 \bitmapfigure[max width=0.44\linewidth]{img/2015/shaanxi_liberal/q20_fig1.png}
 \end{problem}

\begin{solution}
\begin{enumerate}
\item
由点 \(A(0,-1)\) 在椭圆上，得 \(\dfrac{1}{b^2}=1\)，所以 \(b=1\). 又离心率 \(e=\dfrac{c}{a}=\dfrac{\sqrt{2}}{2}\)，而 \(c^2=a^2-b^2\)，故
\[
\frac{a^2-1}{a^2}=\frac{1}{2}
\]
解得 \(a^2=2\). 因此椭圆 \(E\) 的方程为 \(\dfrac{x^2}{2}+y^2=1\).

\item
设直线 \(PQ\) 的方程为 \(y=kx+1-k\). 将其代入椭圆方程，得
\[
\left(k^2+\frac{1}{2}\right)x^2+2k(1-k)x+k(k-2)=0
\]
设两个交点的横坐标为 \(x_1\)，\(x_2\). 若某个 \(x_i=0\)，由椭圆方程可知对应点只能是 \((0,1)\) 或 \(A=(0,-1)\). 后者与条件矛盾；前者代入直线方程得 \(k=0\)，此时上式化为 \(\frac{1}{2}x^2=0\)，两个交点重合，也与 \(P\)，\(Q\) 为不同两点矛盾. 因此 \(x_1x_2\ne0\)，由二次方程的常数项还知 \(k\ne2\)，故由韦达定理，
\[
\frac{1}{x_1}+\frac{1}{x_2}=\frac{x_1+x_2}{x_1x_2}=\frac{2(1-k)}{2-k}
\]
直线 \(AP\)，\(AQ\) 的斜率分别为
\(m_1=\frac{kx_1+1-k+1}{x_1}=k+\frac{2-k}{x_1}\)，\(m_2=k+\frac{2-k}{x_2}\).
所以
\[
m_1+m_2=2k+(2-k)\left(\frac{1}{x_1}+\frac{1}{x_2}\right)=2k+2(1-k)=2
\]
故直线 \(AP\) 与 \(AQ\) 的斜率之和为 \(2\).
\end{enumerate}
\end{solution}

 \begin{problem}
 设 \(f_{n}(x)=x+x^{2}+\cdots+x^{n}-1\)，\(x \geqslant 0\)，\(n \in \N\)，\(n \geqslant 2\).

\begin{enumerate}
\item 求 \(f_{n}^{\prime}(2)\)；
\item 证明：\(f_{n}(x)\) 在 \(\left(0, \frac{2}{3}\right)\) 内有且仅有一个零点（记为 \(a_{n}\)），且 \(0 < a_{n} - \frac{1}{2} < \frac{1}{3}\left(\frac{2}{3}\right)^{n}\).
 \end{enumerate}
 \end{problem}

\begin{solution}
\begin{enumerate}
\item
由
\[
f_n'(x)=1+2x+3x^2+\cdots+nx^{n-1}
\]
并利用 \(x+x^2+\cdots+x^n=\dfrac{x^{n+1}-x}{x-1}\)，得
\[
f_n'(2)=\left.\frac{\mathrm{d}}{\mathrm{d}x}\left(\frac{x^{n+1}-x}{x-1}\right)\right|_{x=2}=(n-1)2^n+1
\]

\item
当 \(x\geqslant0\) 时，\(f_n'(x)>0\)，所以 \(f_n(x)\) 在 \([0,+\infty)\) 上严格递增. 又
\(f_n(0)=-1<0\)，\(f_n\left(\frac{2}{3}\right)=\sum_{j=1}^{n}\left(\frac{2}{3}\right)^j-1=1-2\left(\frac{2}{3}\right)^n>0\)
其中最后一个不等式由 \(n\geqslant2\) 得到. 因此由零点存在性定理，\(f_n(x)\) 在 \(\left(0,\frac{2}{3}\right)\) 内至少有一个零点；再由严格递增性可知该零点唯一，记为 \(a_n\).

计算
\[
f_n\left(\frac{1}{2}\right)=\sum_{j=1}^{n}\left(\frac{1}{2}\right)^j-1=-\left(\frac{1}{2}\right)^n<0
\]
故 \(a_n>\frac{1}{2}\). 令 \(x_0=\frac{1}{2}+\frac{1}{3}\left(\frac{2}{3}\right)^n\). 由 \(n\geqslant2\) 知 \(\frac{1}{3}\left(\frac{2}{3}\right)^n\leqslant\frac{4}{27}<\frac{1}{6}\)，所以 \(x_0<\frac{2}{3}\)，且
\(f_n(x)=\frac{2x-1-x^{n+1}}{1-x}\)，\((x\ne1)\).
由于
\[
2x_0-1=\frac{2}{3}\left(\frac{2}{3}\right)^n=\left(\frac{2}{3}\right)^{n+1}>x_0^{n+1}
\]
所以 \(f_n(x_0)>0\). 由 \(f_n\) 的严格递增性，得 \(a_n<x_0\)，即
\[
0<a_n-\frac{1}{2}<\frac{1}{3}\left(\frac{2}{3}\right)^n
\]
\end{enumerate}
\end{solution}

 \begin{problem}
 如图，\(AB\) 切 \(\odot O\) 于点 \(B\)，直线 \(AO\) 交 \(\odot O\) 于 \(D\)，\(E\) 两点，\(BC \perp DE\)，垂足为 \(C\).

\begin{enumerate}
\item 证明：\(\angle CBD = \angle DBA\)；
\item 若 \(AD = 3DC\)，\(BC = \sqrt{2}\)，求 \(\odot O\) 的直径.
\end{enumerate}

 \bitmapfigure[max width=0.48\linewidth]{img/2015/shaanxi_liberal/q22_fig1.png}
 \end{problem}

\begin{solution}
\begin{enumerate}
\item
因为 \(AO\) 经过圆心，\(DE\) 是圆 \(O\) 的直径，所以 \(\angle DBE=90^\circ\). 由切线弦定理，\(\angle DBA=\angle DEB\). 在直角三角形 \(BDE\) 中，
\[
\angle DBA=\angle DEB=90^\circ-\angle BDE
\]
又因为 \(D\)，\(C\)，\(E\) 共线，且 \(BC\perp DE\)，所以 \(\angle BCD=90^\circ\)，并且 \(\angle BDE=\angle BDC\). 因此在直角三角形 \(BCD\) 中，
\[
\angle CBD=90^\circ-\angle BDC=\angle DBA
\]

\item
设 \(DC=x>0\)，圆的半径为 \(R\). 由图中点的顺序，\(OC=R-x\). 在直角三角形 \(OCB\) 中，\(BC=\sqrt{2}\)，故
\[
R^2=(R-x)^2+2,
\qquad\text{即}\qquad 2Rx=x^2+2
\]
又 \(A\)，\(D\)，\(C\) 共线，且 \(AD=3DC=3x\)，所以 \(AC=4x\). 在直角三角形 \(ABC\) 中，
\[
AB^2=AC^2+BC^2=16x^2+2
\]
由切线与割线定理，\(AB^2=AD\cdot AE=3x(3x+2R)=9x^2+6Rx\)，从而
\[
7x^2+2=6Rx
\]
结合 \(2Rx=x^2+2\)，得 \(7x^2+2=3x^2+6\)，所以 \(x=1\). 再由 \(2Rx=x^2+2\)，得 \(2R=3\). 故圆 \(O\) 的直径为 \(2R=3\).
\end{enumerate}
\end{solution}

 \begin{problem}
 在直角坐标系 \(xOy\) 中，直线 \(l\) 的参数方程为 \(\begin{cases}
x = 3 + \frac{1}{2} t,\\
y = \frac{\sqrt{3}}{2} t,
\end{cases}\)（\(t\) 为参数），以原点为极点，\(x\) 轴正半轴为极轴建立极坐标系. \(\odot C\) 的极坐标方程为 \(\rho = 2\sqrt{3}\sin \theta\).

\begin{enumerate}
\item 写出 \(\odot C\) 的直角坐标方程；
\item \(P\) 为直线 \(l\) 上一动点，当 \(P\) 到圆心 \(C\) 的距离最小时，求 \(P\) 的直角坐标.
 \end{enumerate}
 \end{problem}

\begin{solution}
\begin{enumerate}
\item
由极坐标与直角坐标的关系 \(x=\rho\cos\theta\)，\(y=\rho\sin\theta\)，得
\[
x^2+y^2=\rho^2=2\sqrt{3}\rho\sin\theta=2\sqrt{3}y
\]
所以圆 \(C\) 的直角坐标方程为
\[
x^2+(y-\sqrt{3})^2=3
\]

\item
由直线的参数方程，点 \(P\) 的坐标为 \(P\left(3+\frac{t}{2},\frac{\sqrt{3}}{2}t\right)\)，圆心为 \(C=(0,\sqrt{3})\). 于是
\[
PC^2=\left(3+\frac{t}{2}\right)^2+\left(\frac{\sqrt{3}}{2}t-\sqrt{3}\right)^2=t^2+12\geqslant12
\]
当且仅当 \(t=0\) 时取等号，此时 \(P=(3,0)\).
\end{enumerate}
\end{solution}

 \begin{problem}
 已知关于 \(x\) 的不等式 \(|x + a| < b\) 的解集为 \(\{x \mid 2 < x < 4\}\).

\begin{enumerate}
\item 求实数 \(a\)，\(b\) 的值；
\item 求 \(\sqrt{at+12}+\sqrt{bt}\) 的最大值.
 \end{enumerate}
 \end{problem}

\begin{solution}
\begin{enumerate}
\item
由
\[
|x+a|<b\Longleftrightarrow -b<x+a<b\Longleftrightarrow -a-b<x<-a+b
\]
题设解集为 \(2<x<4\)，故
\(-a-b=2\)，\(-a+b=4\).
两式相加、相减，得 \(a=-3\)，\(b=1\).

\item
因此根式有意义时 \(0\leqslant t\leqslant4\)，且
\[
\sqrt{at+12}+\sqrt{bt}=\sqrt{12-3t}+\sqrt{t}=\sqrt{3}\sqrt{4-t}+\sqrt{t}
\]
由柯西不等式，
\[
\sqrt{3}\sqrt{4-t}+\sqrt{t}\leqslant\sqrt{3+1}\sqrt{(4-t)+t}=4
\]
当 \(\sqrt{4-t}:\sqrt{t}=\sqrt{3}:1\)，即 \(t=1\) 时取等号. 因此最大值为 \(4\).
\end{enumerate}
\end{solution}
